\section{Convergence: CRDT, CALM, and predicatory confluence}

This section records claim nodes and dependency edges.

\begin{proposition}[Confluence (order form)]\label{thm:confluent}\lean{Infomaterial.CRDT.Confluent}\leanok\uses{thm:settling, thm:valuation, thm:four-bilattice}
Two valuations $u$, $w$ are \emph{confluent} when they have a common upper bound in the knowledge order: $\exists z,\ u \leq_k z \wedge w \leq_k z$. The divergent settling paths from a common source reach a common $z$ by adding evidence --- predicatory confluence in its order reading.
\end{proposition}
\begin{theorem}[Confluence is total --- the join is the witness]\label{thm:confluent-total}\lean{Infomaterial.CRDT.confluent_joinK}\leanok\uses{thm:confluent, thm:section-lub, thm:valuation}
For \emph{any} two valuations the knowledge join $u \vee_k w$ is a common upper bound, so \texttt{Confluent u w} always holds (\texttt{confluent\_joinK}, reusing the lub laws of \cref{thm:section-lub}). Predicatory confluence is therefore not a contingent fact about particular reductions but a structural fact of the knowledge order: the least upper bound always exists, and operational posit/merge paths reach it because they are $\leq_k$-increasing evidence additions.
\end{theorem}
\begin{theorem}[CRDT convergence --- order-independent merge]\label{thm:crdt-order-independent}\lean{Infomaterial.CRDT.crdt_merge_order_independent}\leanok\uses{thm:vanishing-h1-crdt, thm:charge-joinList}
The knowledge join of a list of deltas depends only on the set of deltas, not their order or multiplicity (\texttt{crdt\_merge\_order\_independent}, a re-export of \cref{thm:vanishing-h1-crdt}). Replicas integrating the same set of deltas converge to the same value without coordination --- ``deltas integrate without coordination'' $=$ $H^1 = 0$ in the finite regime.
\end{theorem}
\begin{theorem}[Convergence facts of the confluent modal settling]\label{thm:calm-convergence-facts}\lean{Infomaterial.CRDT.calm_convergence_facts}\leanok\uses{thm:h1-vanishes, thm:vanishing-h1-crdt, thm:confluent-total, thm:charge-joinList, thm:valuation, thm:settling}
For a modal settling frame, the four convergence readings all \emph{hold} (\texttt{calm\_convergence\_facts}): \textbf{(i)} $H^1 = 0$ --- every agreeing family glues to a unique global section (\texttt{h1\_vanishes}); \textbf{(ii)} order-independent merge (\cref{thm:vanishing-h1-crdt}); \textbf{(iii)} confluence --- every pair joins at $u \vee_k w$; \textbf{(iv)} unique settled fixed point (quiescence $=$ closure fixed point, with $\leq_k$ antisymmetric). This is a \emph{conjunction}: each holds \emph{unconditionally} in this lattice model (confluence is total because $\vee_k$ always exists; $H^1 = 0$ because every cocycle glues), so no implication \emph{between} the readings is established. The genuine CALM biconditional is \cref{thm:calm-convergence-equivalence}.
\end{theorem}
\begin{theorem}[CALM convergence equivalence]\label{thm:calm-convergence-equivalence}\notready\uses{thm:h1-vanishes, thm:confluent}
\emph{(Open frontier.)} The substantive CALM theorem: $H^1 = 0 \iff$ confluence $\iff$ monotone, as a genuine \emph{biconditional} that distinguishes the coordination-free case from a non-monotone settling with $H^1 \neq 0$. The confluent-lattice model formalized above makes both confluence and $H^1$-vanishing unconditional, so it cannot witness the equivalence (there is no non-confluent / $H^1 \neq 0$ side to separate); a model with genuine operational (posit/merge-path) confluence and non-vanishing cohomology is required. Red until such a model and the biconditional are constructed.
\end{theorem}
\begin{theorem}[CALM --- monotone settling is coordination-free]\label{thm:calm-monotone}\lean{Infomaterial.CRDT.calm_monotone_coordination_free}\leanok\uses{thm:vanishing-h1-crdt, thm:charge-joinList, thm:valuation, thm:predication-seed, thm:settling}
For a modal settling frame and any seed, the closure orbit is $\leq_k$-monotone in the settle index --- no information is retracted across steps (the CALM monotonicity hypothesis) --- and the delta merge is order/multiplicity independent (\cref{thm:vanishing-h1-crdt}). Bundled (\texttt{calm\_monotone\_coordination\_free}), this is the CALM theorem's projection on settling: a monotone computation can be merged in any order with the same result, so replicas need no coordination.
\end{theorem}
\begin{theorem}[Mayer--Vietoris overlap gluing]\label{thm:mayer-vietoris}\lean{Infomaterial.CRDT.mayer_vietoris_overlap_glues}\leanok\uses{thm:section-lub, thm:sheaf-cover}
The pre-images of a closure-stable value $w$ are closed under the overlap join: if $\operatorname{cl} v = w$ and $\operatorname{cl} v' = w$ then $\operatorname{cl}(v \vee_k v') = w$ (\texttt{mayer\_vietoris\_overlap\_glues}). This cosettling join-subsemilattice is the algebraic projection of Mayer--Vietoris gluing: local cosettling data on a cover glue along their overlaps, the open-cover structure behind $H^1 = 0$.
\end{theorem}
